\section{Conclusion}
\label{sec-conclusion}

Optimizing LLM inference for high throughput and low latency is desirable but challenging. We presented a broad characterization of existing LLM inference schedulers by dividing them into two categories -- \prefillprio and \decodeprio. In general, we argue that the former category is better at optimizing throughput whereas the latter is better at optimizing TBT latency. However, none of them is ideal when optimizing throughput and latency are both important. 

To address this tradeoff, we introduce \sysname --- a system that instantiates a novel approach comprised of  \chunking and \hbatch. \sysname chunks input prompts into smaller units of work to create stall-free schedules. This way, \sysname can add new requests in a running batch without pausing ongoing decodes. Our evaluation shows that \sysname improves the serving capacity of \mistral by up to 2.6\myx on a single A100 GPU and up to 5.6\myx for \falcon on 8 A100 GPUs.




\section{Acknowledgement}
We would like to thank OSDI reviewers and our shepherd for their insightful feedback. This research is partly supported by GT Cloud Hub, under the auspices of the Institute
for Data Engineering and Science (IDEaS), with funding from Microsoft, and the Center for Research into Novel Compute Hierarchies (CRNCH) at Georgia Tech.
